% `tex.web` §963 makes the language the first character of every trie key:
%
%     q:=0; hc[0]:=cur_lang;
%     while l<=k do …
%
% so patterns loaded under one `\language` are invisible under another, and so
% are `\hyphenation` exceptions (§934 keys those the same way).  §1034 says
% what `cur_lang` is:
%
%     if language<=0 then cur_lang:=0
%     else if language>255 then cur_lang:=0
%     else cur_lang:=language
%
% — out of range is language nought, not an error.
%
% We kept one set for every language, so patterns loaded under `\language5`
% hyphenated a paragraph set under `\language0`.  A document that loads several
% languages' patterns — which is what every multilingual format does — had all
% of them applying at once.
\catcode`\{=1 \catcode`\}=2
\font\r=cmr10 \r \hsize20pt \parindent0pt \baselineskip12pt
\tolerance10000 \pretolerance-1 \hbadness10000 \uchyph=1
\showboxbreadth99 \showboxdepth99 \tracingonline0
\language=5 \patterns{bb1}
\language=0 \patterns{cc1}
% Under language 0: the `cc` patterns apply and the `bb` ones do not.
\setbox0\vbox{\noindent x BBBBBB\par}\showbox0
\setbox1\vbox{\noindent x CCCCCC\par}\showbox1
% Under language 5: the other way about.
\language=5
\setbox2\vbox{\noindent x BBBBBB\par}\showbox2
\setbox3\vbox{\noindent x CCCCCC\par}\showbox3
% Out of range is language nought.
\language=300
\setbox4\vbox{\noindent x CCCCCC\par}\showbox4
% A language assignment inside a paragraph becomes a language whatsit before
% the next character; the breaker changes pattern tries at that exact point.
\language=0
\setbox5\vbox{\noindent x BBBBBB \language=5 CCCCCC\par}\showbox5
% Explicit `\setlanguage` changes the same breaker state without assigning the
% `\language` parameter.
\setbox6\vbox{\noindent x \setlanguage5 CCCCCC\par}\showbox6
\end
